\documentclass[aps,preprint]{revtex4}%
\usepackage{amsfonts}
\usepackage{amsmath}
\usepackage{amssymb}
\usepackage{graphicx}%
\setcounter{MaxMatrixCols}{30}
%TCIDATA{OutputFilter=latex2.dll}
%TCIDATA{Version=4.10.0.2363}
%TCIDATA{CSTFile=revtex4.cst}
%TCIDATA{Created=Monday, July 21, 2003 20:57:35}
%TCIDATA{LastRevised=Thursday, July 24, 2003 22:22:27}
%TCIDATA{<META NAME="GraphicsSave" CONTENT="32">}
%TCIDATA{<META NAME="DocumentShell" CONTENT="Articles\SW\REVTeX 4">}
\newtheorem{theorem}{Theorem}
\newtheorem{acknowledgement}[theorem]{Acknowledgement}
\newtheorem{algorithm}[theorem]{Algorithm}
\newtheorem{axiom}[theorem]{Axiom}
\newtheorem{claim}[theorem]{Claim}
\newtheorem{conclusion}[theorem]{Conclusion}
\newtheorem{condition}[theorem]{Condition}
\newtheorem{conjecture}[theorem]{Conjecture}
\newtheorem{corollary}[theorem]{Corollary}
\newtheorem{criterion}[theorem]{Criterion}
\newtheorem{definition}[theorem]{Definition}
\newtheorem{example}[theorem]{Example}
\newtheorem{exercise}[theorem]{Exercise}
\newtheorem{lemma}[theorem]{Lemma}
\newtheorem{notation}[theorem]{Notation}
\newtheorem{problem}[theorem]{Problem}
\newtheorem{proposition}[theorem]{Proposition}
\newtheorem{remark}[theorem]{Remark}
\newtheorem{solution}[theorem]{Solution}
\newtheorem{summary}[theorem]{Summary}
\newenvironment{proof}[1][Proof]{\noindent\textbf{#1.} }{\ \rule{0.5em}{0.5em}}
\begin{document}
\preprint{HEP/123-qed}
\title[Short title for running header]{Long title that appears on the title page}
\author{First Author}
\affiliation{My Institution}
\author{Second Author}
\affiliation{My Institution}
\author{Third Author}
\affiliation{Other Institution}
\keywords{one two three}
\pacs{PACS number}

\begin{abstract}
Shell document for REV\TeX{} 4.

\end{abstract}
\volumeyear{year}
\volumenumber{number}
\issuenumber{number}
\eid{identifier}
\date[Date text]{date}
\received[Received text]{date}

\revised[Revised text]{date}

\accepted[Accepted text]{date}

\published[Published text]{date}

\startpage{101}
\endpage{102}
\maketitle
\tableofcontents


\section{Occupation Number Functions}

\subsection{First Order Reduced Density Operator in Terms of Geminal
Occupation Numbers}%

\begin{equation}
N_{j}=\frac{n_{j}S_{N-1}\left(  \widehat{n_{j}}\right)  }{S_{N}\left(
\mathbf{n}\right)  }%
\end{equation}
defines a positive map
\begin{equation}
\mathbf{N}:\mathbb{R}^{s}\rightarrow\mathbb{R}^{s}%
\end{equation}
such that
\begin{equation}
\mathbf{N=}\left(
\begin{array}
[c]{c}%
N_{1}\\
\vdots\\
N_{s}%
\end{array}
\right)  =\left(
\begin{array}
[c]{c}%
N_{1}\left(  \mathbf{n}\right) \\
\vdots\\
N_{s}\left(  \mathbf{n}\right)
\end{array}
\right)  =\frac{1}{S_{N}\left(  \mathbf{n}\right)  }\left(
\begin{array}
[c]{c}%
n_{1}S_{N-1}\left(  \widehat{n_{1}}\right) \\
\vdots\\
n_{s}S_{N-1}\left(  \widehat{n_{s}}\right)
\end{array}
\right)
\end{equation}
where
\begin{equation}
S_{N}\left(  \mathbf{n}\right)  =\sum_{1\leq j_{1}<\cdots j_{s}\leq s}%
n_{j_{1}}\cdots n_{j_{N}}%
\end{equation}
and
\begin{equation}
S_{N-q}\left(  \widehat{n_{1}}\cdots\widehat{n_{q}}\right)  =\sum
_{\substack{1\leq j_{1}<\cdots j_{N-q}\leq s\\j_{1},\cdots,j_{N-q}\neq
k_{1},\cdots,k_{q}}}n_{j_{1}}\cdots n_{j_{N-q}}%
\end{equation}
One should observe that if $0\leq n_{j}\leq1$ then $0\leq N_{j}\leq1;1\leq
j\leq s$. The functions $\left\{  N_{j}\left(  \mathbf{n}\right)  \right\}  $
are analytic and have Taylor series expansion
\begin{align}
N_{j}\left(  \mathbf{n}_{0}+\delta\mathbf{n}\right)   &  =\sum_{k=0}^{\infty
}\frac{1}{k!}\left\langle D^{k}N_{j}\left(  \mathbf{n}_{0}\right)  \left\vert
\overset{k}{\overbrace{\delta\mathbf{n}\otimes\cdots\otimes\delta\mathbf{n}}%
}\right.  \right\rangle \\
&  =\sum_{k=0}^{\infty}\sum_{1\leq j_{1},\ldots,j_{k}\leq s}\frac{1}{k!}%
D^{k}N_{j}\left(  \mathbf{n}_{0}\right)  _{j_{1}\cdots j_{k}}\delta n_{j_{1}%
}\cdots\delta n_{j_{k}}%
\end{align}
where
\begin{equation}
\left[  D^{k}N_{j}\left(  \mathbf{n}_{0}\right)  \right]  _{j_{1}\cdots j_{k}%
}=\frac{\partial^{k}N_{j}\left(  \mathbf{n}_{0}\right)  }{\partial n_{j_{1}%
}\cdots\partial n_{j_{k}}};\quad1\leq j_{1},\ldots,j_{k}\leq s
\end{equation}%
\begin{equation}
\delta\mathbf{n}=\mathbf{n}-\mathbf{n}_{0}%
\end{equation}
and
\begin{equation}
D^{0}N_{j}\left(  \mathbf{n}_{0}\right)  =N_{j}\left(  \mathbf{n}_{0}\right)
\end{equation}
The full vector map is then
\begin{equation}
\mathbf{N}\left(  \mathbf{n}_{0}+\delta\mathbf{n}\right)  =\sum_{k=0}^{\infty
}\frac{1}{k!}D^{k}\mathbf{N}\left(  \mathbf{n}_{0}\right)  \left(
\delta\mathbf{n}\otimes\cdots\otimes\delta\mathbf{n}\right)
\end{equation}
In particular
\begin{align}
\frac{\partial N_{j}\left(  \mathbf{n}_{0}\right)  }{\partial n_{l}}  &
=\frac{\partial}{\partial n_{l}}\left\{  \frac{n_{j}S_{N}\left(
\widehat{n_{j}}\right)  }{S_{N}\left(  \mathbf{n}\right)  }\right\}
\nonumber\\
&  =\left.  -\frac{1}{S_{N}\left(  \mathbf{n}\right)  ^{2}}\frac{\partial
S_{N}}{\partial n_{l}}n_{j}S_{N}\left(  \widehat{n_{j}}\right)  +\frac
{1}{S_{N}\left(  \mathbf{n}\right)  }S_{N}\left(  \widehat{n_{j}}\right)
\delta_{jl}+\frac{1}{S_{N}\left(  \mathbf{n}\right)  }n_{j}\frac{\partial
S_{N}\left(  \widehat{n_{j}}\right)  }{\partial n_{l}}\right\vert
_{\mathbf{n=n}_{0}}\nonumber\\
&  =\left.  \frac{1}{S_{N}\left(  \mathbf{n}\right)  }\left\{  S_{N}\left(
\widehat{n_{j}}\right)  \delta_{jl}+n_{j}S_{N}\left(  \widehat{n_{j}}%
\widehat{n_{l}}\right)  -n_{j}S_{N}\left(  \widehat{n_{l}}\right)
S_{N}\left(  \widehat{n_{j}}\right)  \right\}  \right\vert _{\mathbf{n=n}_{0}}%
\end{align}
One can make the substitution
\begin{equation}
n_{j}=e^{\eta_{j}}\Leftrightarrow\ln n_{j}=\eta_{j}\quad-\infty\leq\eta
_{j}\leq\infty
\end{equation}
then
\begin{equation}
N_{j}=\frac{\mathcal{S}_{N-1}\left(  \widehat{\eta_{j}}\right)  }%
{\mathcal{S}_{N}\left(  \mathbf{\eta}\right)  }=\frac{\partial\ln
\mathcal{S}_{N}\left(  \mathbf{\eta}\right)  }{\partial\eta_{j}}%
\end{equation}
where
\begin{align}
\mathcal{S}_{N}\left(  \mathbf{\eta}\right)   &  =\sum_{1\leq j_{1}<\cdots
j_{s}\leq s}e^{\eta_{j_{1}}+\cdots+\eta_{j_{N}}}\\
\mathcal{S}_{N-q}\left(  \widehat{\eta_{1}}\cdots\widehat{\eta_{q}}\right)
&  =\sum_{\substack{1\leq j_{1}<\cdots j_{N-q}\leq s\\j_{1},\cdots,j_{N-q}\neq
k_{1},\cdots,k_{q}}}e^{\eta_{j_{1}}+\cdots+\eta_{j_{N-q}}}%
\end{align}
We then obtain the Taylor expansion
\begin{equation}
\mathbf{N}\left(  \mathbf{\eta}_{0}+\delta\mathbf{\eta}\right)  =\sum
_{k=0}^{\infty}\frac{1}{k!}D^{k}\mathbf{N}\left(  \mathbf{\eta}_{0}\right)
\left(  \delta\mathbf{\eta}\otimes\cdots\otimes\delta\mathbf{\eta}\right)
\end{equation}
where
\begin{equation}
\delta\mathbf{\eta}=\mathbf{\eta}-\mathbf{\eta}_{0}%
\end{equation}
A natural point to expand about would be
\begin{equation}
\mathbf{n}_{0}=\left(
\begin{array}
[c]{c}%
\left.
\begin{array}
[c]{c}%
1\\
\vdots\\
1
\end{array}
\right\}  N\\
\qquad\left.
\begin{array}
[c]{c}%
0\\
\vdots\\
0
\end{array}
\right\}  s-N
\end{array}
\right)  \leftrightarrow\mathbf{\eta}_{0}=\left(
\begin{array}
[c]{c}%
\left.
\begin{array}
[c]{c}%
0\\
\vdots\\
0
\end{array}
\right\}  N\\
\qquad\left.
\begin{array}
[c]{c}%
-\infty\\
\vdots\\
-\infty
\end{array}
\right\}  s-N
\end{array}
\right)
\end{equation}
which corresponds to a $2N$-particle IPS, however numerically this is not very
satisfactory. Another point could be (as long as $s<\infty$)
\begin{equation}
\mathbf{n}_{0}=\left(  \left.
\begin{array}
[c]{c}%
1\\
\vdots\\
1
\end{array}
\right\}  s\right)  \leftrightarrow\mathbf{\eta}_{0}=\left(  \left.
\begin{array}
[c]{c}%
0\\
\vdots\\
0
\end{array}
\right\}  s\right)
\end{equation}


The first derrivatives are then
\begin{align}
\frac{\partial N_{j}\left(  \mathbf{\eta}\right)  }{\partial\eta_{l}} &
=\frac{\partial}{\partial\eta_{l}}\left\{  \frac{\partial\ln\mathcal{S}%
_{N}\left(  \mathbf{\eta}\right)  }{\partial\eta_{j}}\right\}  =\frac
{1}{\mathcal{S}_{N}\left(  \mathbf{\eta}\right)  }\left\{  \frac
{\partial\mathcal{S}_{N-1}\left(  \widehat{\eta_{j}}\right)  }{\partial
\eta_{l}}-\frac{1}{\mathcal{S}_{N}\left(  \mathbf{\eta}\right)  }%
\frac{\partial\mathcal{S}_{N}\left(  \mathbf{\eta}\right)  }{\partial\eta_{l}%
}\frac{\partial\mathcal{S}_{N}\left(  \mathbf{\eta}\right)  }{\partial\eta
_{j}}\right\}  \nonumber\\
&  =\frac{1}{\mathcal{S}_{N}\left(  \mathbf{\eta}\right)  }\left\{
\frac{\partial^{2}\mathcal{S}_{N}\left(  \widehat{\eta_{j}}\right)  }%
{\partial\eta_{l}\partial\eta_{j}}-\frac{1}{\mathcal{S}_{N}\left(
\mathbf{\eta}\right)  }\frac{\partial\mathcal{S}_{N}\left(  \mathbf{\eta
}\right)  }{\partial\eta_{l}}\frac{\partial\mathcal{S}_{N}\left(
\mathbf{\eta}\right)  }{\partial\eta_{j}}\right\}  \nonumber\\
&  =\frac{1}{\mathcal{S}_{N}\left(  \mathbf{\eta}\right)  }\left\{
\mathcal{S}_{N-2}\left(  \widehat{\eta_{j}}\widehat{\eta_{l}}\right)
-\frac{1}{\mathcal{S}_{N}\left(  \mathbf{\eta}\right)  }\mathcal{S}%
_{N-1}\left(  \widehat{\eta_{j}}\right)  \mathcal{S}_{N-1}\left(
\widehat{\eta_{l}}\right)  \right\}
\end{align}
This can also be expressed in vector notation as
\begin{equation}
D^{1}\mathbf{N}\left(  \mathbf{\eta}\right)  =\frac{1}{\mathcal{S}_{N}\left(
\mathbf{\eta}\right)  }\left\{  D^{2}\mathcal{S}_{N}\left(  \mathbf{\eta
}\right)  -\frac{1}{\mathcal{S}_{N}\left(  \mathbf{\eta}\right)  }%
D^{1}\mathcal{S}_{N}\left(  \mathbf{\eta}\right)  \otimes D^{1}\mathcal{S}%
_{N}\left(  \mathbf{\eta}\right)  \right\}
\end{equation}
If one evaluates $D^{1}\mathbf{N}\left(  \mathbf{\eta}\right)  $ at
$\mathbf{\eta}=0$ $\leftrightarrow\mathbf{n}=1$ then
\begin{equation}
\frac{\partial N_{j}\left(  \mathbf{0}\right)  }{\partial\eta_{l}}%
=\frac{\partial N_{j}\left(  \mathbf{1}\right)  }{\partial n_{l}}=\left.
\frac{1}{S_{N}\left(  \mathbf{n}\right)  }\left\{  S_{N}\left(  \widehat
{n_{j}}\right)  \delta_{jl}+n_{j}S_{N}\left(  \widehat{n_{j}}\widehat{n_{l}%
}\right)  -n_{j}S_{N}\left(  \widehat{n_{l}}\right)  S_{N}\left(
\widehat{n_{j}}\right)  \right\}  \right\vert _{\mathbf{n=1}}%
\end{equation}
which gives
\begin{align}
S_{N}\left(  \mathbf{1}\right)   &  =\binom{s}{N}\\
S_{N}\left(  \widehat{n_{j}}\right)   &  =\binom{s-1}{N-1}\\
S_{N}\left(  \widehat{n_{j}}\widehat{n_{l}}\right)   &  =\binom{s-2}{N-2}%
\end{align}
hence%
\begin{align}
\frac{\partial N_{j}\left(  \mathbf{0}\right)  }{\partial\eta_{l}} &
=\frac{N!\left(  s-N\right)  !}{s!}\left\{  \frac{\left(  s-1\right)
!}{\left(  N-1\right)  !\left(  s-N\right)  !}\delta_{jl}+\frac{\left(
s-2\right)  !}{\left(  N-2\right)  !\left(  s-N\right)  !}-\left(
\frac{\left(  s-1\right)  !}{\left(  N-1\right)  !\left(  s-N\right)
!}\right)  ^{2}\right\}  \\
&  =\frac{N}{s}\delta_{jl}-\frac{N\left(  N-1\right)  }{s\left(  s-1\right)
}-\frac{\left(  s-1\right)  !}{\left(  N-1\right)  !\left(  s-N\right)
!}\frac{N}{s}%
\end{align}


\subsection{Fixed AGP Normalization}

Let us consider the hypersurface
\begin{equation}
S_{N}\left(  \mathbf{\eta}\right)  =1
\end{equation}
then
\begin{equation}
\mathbf{N}\left(  \mathbf{\eta}\right)  =\nabla S_{N}\left(  \mathbf{\eta
}\right)
\end{equation}
and
\begin{equation}
\mathbf{N}\left(  \mathbf{\eta}\right)  ^{t}\cdot\delta\mathbf{\eta}=0
\end{equation}
where $\delta\mathbf{\eta\in}T_{S_{N}\left(  \mathbf{\eta}\right)  }\equiv$
the tangent space of $S_{N}\left(  \mathbf{\eta}\right)  -1$ at the point
$\mathbf{\eta}$, i.e $\mathbf{N}\left(  \mathbf{\eta}\right)  $ is the normal
vector to $T_{S_{N}\left(  \mathbf{\eta}\right)  }^{\prime}$ 

\subsection{Geminal Occupation Numbers in Terms of First Order Reduced Density
Operator Occupation Numbers (The Inverse Map)}

The Geminal Occupation Number maps are also a positive map%

\begin{align}
\mathbf{n}  &  :\mathbb{R}^{s}\rightarrow\mathbb{R}^{s}\\
\mathbf{\eta}  &  :\mathbb{R}^{s}\rightarrow\mathbb{R}^{s}%
\end{align}
where
\begin{align}
\mathbf{n}\left(  \mathbf{N}\left(  \mathbf{n}\right)  \right)   &  =\left(
\mathbf{n\circ N}\right)  \left(  \mathbf{n}\right)  =\mathbf{n}\\
\mathbf{\eta}\left(  \mathbf{N}\left(  \mathbf{\eta}\right)  \right)   &
=\left(  \mathbf{\eta\circ N}\right)  \left(  \mathbf{\eta}\right)
=\mathbf{\eta}%
\end{align}
so that both
\begin{align}
\mathbf{n}  &  =\mathbf{N}^{-1}\\
\mathbf{\eta}  &  =\mathbf{N}^{-1}%
\end{align}
but in different coordinate systems. The inverse map has Taylor series
expansion
\begin{equation}
\mathbf{\eta}\left(  \mathbf{N}_{0}+\delta\mathbf{N}\right)  =\sum
_{k=0}^{\infty}\frac{1}{k!}D^{k}\mathbf{\eta}\left(  \mathbf{N}_{0}\right)
\left(  \delta\mathbf{N}\otimes\cdots\otimes\delta\mathbf{N}\right)
\end{equation}
The composition map, which is the identity map, has a Taylor series expansion
\begin{equation}
\left(  \mathbf{\eta\circ N}\right)  \left(  \mathbf{\eta}_{0}+\delta
\mathbf{\eta}\right)  =\mathbf{\eta}_{0}+\delta\mathbf{\eta}%
\end{equation}
which implies that
\begin{equation}
D^{k}\left(  \mathbf{\eta\circ N}\right)  \left(  \mathbf{\eta}_{0}\right)
=\mathbf{0,\quad}k>1
\end{equation}
and
\begin{equation}
D^{1}\left(  \mathbf{\eta\circ N}\right)  \left(  \mathbf{\eta}_{0}\right)
=D^{1}\mathbf{\eta}\left(  \mathbf{N}_{0}\right)  \mathbf{\circ}%
D^{1}\mathbf{N}\left(  \mathbf{\eta}_{0}\right)  =\mathbf{I}%
\end{equation}
where
\begin{align}
\mathbf{N}_{0}  &  =\mathbf{N}\left(  \mathbf{\eta}_{0}\right) \\
\mathbf{\eta}_{0}  &  =\mathbf{\eta}\left(  \mathbf{N}_{0}\right)
\end{align}
Hence%
\begin{equation}
D^{1}\mathbf{\eta}\left(  \mathbf{N}_{0}\right)  =\left[  D^{1}\mathbf{N}%
\left(  \mathbf{\eta}_{0}\right)  \right]  ^{-1}%
\end{equation}
Now
\begin{align}
D^{2}\left(  \mathbf{\eta\circ N}\right)  \left(  \mathbf{\eta}_{0}\right)
&  =D^{1}\left(  D^{1}\left(  \mathbf{\eta\circ N}\right)  \left(
\mathbf{\eta}_{0}\right)  \right)  =D^{1}\left(  D^{1}\mathbf{\eta}\left(
\mathbf{N}_{0}\right)  \mathbf{\circ}D^{1}\mathbf{N}\left(  \mathbf{\eta}%
_{0}\right)  \right) \nonumber\\
&  =D^{2}\mathbf{\eta}\left(  \mathbf{N}_{0}\right)  \mathbf{\circ}%
D^{1}\mathbf{N}\left(  \mathbf{\eta}_{0}\right)  \circ D^{1}\mathbf{N}\left(
\mathbf{\eta}_{0}\right)  +D^{1}\mathbf{\eta}\left(  \mathbf{N}_{0}\right)
\mathbf{\circ}D^{2}\mathbf{N}\left(  \mathbf{\eta}_{0}\right) \nonumber\\
&  =D^{2}\mathbf{\eta}\left(  \mathbf{N}_{0}\right)  \mathbf{\circ}\left(
D^{1}\mathbf{N}\left(  \mathbf{\eta}_{0}\right)  \right)  ^{2}+\left[
D^{1}\mathbf{N}\left(  \mathbf{\eta}_{0}\right)  \right]  ^{-1}\mathbf{\circ
}D^{2}\mathbf{N}\left(  \mathbf{\eta}_{0}\right)  =0
\end{align}%
\begin{equation}
\Rightarrow D^{2}\mathbf{\eta}\left(  \mathbf{N}_{0}\right)  =-\left[
D^{1}\mathbf{N}\left(  \mathbf{\eta}_{0}\right)  \right]  ^{-1}\mathbf{\circ
}D^{2}\mathbf{N}\left(  \mathbf{\eta}_{0}\right)  \circ\left[  D^{1}%
\mathbf{N}\left(  \mathbf{\eta}_{0}\right)  \right]  ^{-2}\qquad\qquad
\qquad\quad
\end{equation}
and
\begin{align}
&  D^{3}\left(  \mathbf{\eta\circ N}\right)  \left(  \mathbf{\eta}_{0}\right)
=D^{1}\left(  D^{2}\mathbf{\eta}\left(  \mathbf{N}_{0}\right)  \mathbf{\circ
}\left[  D^{1}\mathbf{N}\left(  \mathbf{\eta}_{0}\right)  \right]  ^{2}%
+D^{1}\mathbf{\eta}\left(  \mathbf{N}_{0}\right)  \mathbf{\circ}%
D^{2}\mathbf{N}\left(  \mathbf{\eta}_{0}\right)  \right) \nonumber\\
&  =D^{3}\mathbf{\eta}\left(  \mathbf{N}_{0}\right)  \mathbf{\circ}\left[
D^{1}\mathbf{N}\left(  \mathbf{\eta}_{0}\right)  \right]  ^{3}+3D^{2}%
\mathbf{\eta}\left(  \mathbf{N}_{0}\right)  \mathbf{\circ}D^{1}\mathbf{N}%
\left(  \mathbf{\eta}_{0}\right)  \mathbf{\circ}D^{2}\mathbf{N}\left(
\mathbf{\eta}_{0}\right)  +D^{1}\mathbf{\eta}\left(  \mathbf{N}_{0}\right)
\mathbf{\circ}D^{3}\mathbf{N}\left(  \mathbf{\eta}_{0}\right) \nonumber\\
&  =D^{3}\mathbf{\eta}\left(  \mathbf{N}_{0}\right)  \mathbf{\circ}\left[
D^{1}\mathbf{N}\left(  \mathbf{\eta}_{0}\right)  \right]  ^{3}-3\left[
D^{1}\mathbf{N}\left(  \mathbf{\eta}_{0}\right)  \right]  ^{-1}\mathbf{\circ
}D^{2}\mathbf{N}\left(  \mathbf{\eta}_{0}\right)  \circ\left[  D^{1}%
\mathbf{N}\left(  \mathbf{\eta}_{0}\right)  \right]  ^{-1}\mathbf{\circ}%
D^{2}\mathbf{N}\left(  \mathbf{\eta}_{0}\right) \nonumber\\
&  \qquad\qquad\qquad\qquad\qquad\qquad\qquad\qquad\qquad\qquad\qquad
\qquad\qquad+\left[  D^{1}\mathbf{N}\left(  \mathbf{\eta}_{0}\right)  \right]
^{-1}\mathbf{\circ}D^{3}\mathbf{N}\left(  \mathbf{\eta}_{0}\right)  =0
\end{align}%
\begin{align}
&  \Rightarrow D^{3}\mathbf{\eta}\left(  \mathbf{N}_{0}\right)  =\left\{
3\left[  D^{1}\mathbf{N}\left(  \mathbf{\eta}_{0}\right)  \right]
^{-1}\mathbf{\circ}D^{2}\mathbf{N}\left(  \mathbf{\eta}_{0}\right)
\circ\left[  D^{1}\mathbf{N}\left(  \mathbf{\eta}_{0}\right)  \right]
^{-1}\mathbf{\circ}D^{2}\mathbf{N}\left(  \mathbf{\eta}_{0}\right)  \right.
\nonumber\\
&  \qquad\qquad\qquad\qquad\qquad\qquad\qquad\left.  -\left[  D^{1}%
\mathbf{N}\left(  \mathbf{\eta}_{0}\right)  \right]  ^{-1}\mathbf{\circ}%
D^{3}\mathbf{N}\left(  \mathbf{\eta}_{0}\right)  \right\}  \left[
D^{1}\mathbf{N}\left(  \mathbf{\eta}_{0}\right)  \right]  ^{-3}%
\end{align}


\section{Scalar Case}%

\begin{align}
\left(  f\circ g\right)  \left(  x+\delta x\right)   &  =\sum\limits_{k=0}%
^{\infty}\frac{1}{k!}D^{k}\left(  f\circ g\right)  \left(  x\right)  \left(
\delta x\right)  ^{k}=x+\delta x\nonumber\\
&  \Rightarrow D^{k}\left(  f\circ g\right)  \left(  x\right)  =0\quad
k>1\nonumber\\
&  \Rightarrow D^{1}\left(  f\circ g\right)  \left(  x\right)  =1
\end{align}%
\begin{equation}
D^{0}\left(  f\circ g\right)  \left(  x\right)  =\left(  f\circ g\right)
\left(  x\right)
\end{equation}%
\begin{align}
D^{1}\left(  f\circ g\right)  \left(  x\right)   &  =\frac{d}{dx}\left(
f\circ g\right)  \left(  x\right)  =\frac{df\left(  g\right)  }{dg}%
\frac{dg\left(  x\right)  }{dx}\nonumber\\
&  \Rightarrow\frac{df\left(  g\right)  }{dg}=\left(  \frac{dg\left(
x\right)  }{dx}\right)  ^{-1}%
\end{align}
%

\begin{align}
D^{2}\left(  f\circ g\right)  \left(  x\right)   &  =\frac{d}{dx}\left\{
\frac{df\left(  g\right)  }{dg}\frac{dg\left(  x\right)  }{dx}\right\}
=\frac{d^{2}f\left(  g\right)  }{dg^{2}}\left(  \frac{dg\left(  x\right)
}{dx}\right)  ^{2}+\frac{df\left(  g\right)  }{dg}\frac{d^{2}g\left(
x\right)  }{dx^{2}}=0\nonumber\\
&  \Rightarrow\frac{d^{2}f\left(  g\right)  }{dg^{2}}=-\frac{df\left(
g\right)  }{dg}\frac{d^{2}g\left(  x\right)  }{dx^{2}}\left(  \frac{dg\left(
x\right)  }{dx}\right)  ^{-2}=-\left(  \frac{dg\left(  x\right)  }{dx}\right)
^{-1}\frac{d^{2}g\left(  x\right)  }{dx^{2}}\left(  \frac{dg\left(  x\right)
}{dx}\right)  ^{-2}%
\end{align}%
\begin{align}
&  D^{3}\left(  f\circ g\right)  \left(  x\right)  =\frac{d}{dx}\left\{
\frac{d^{2}f\left(  g\right)  }{dg^{2}}\left(  \frac{dg\left(  x\right)  }%
{dx}\right)  ^{2}+\frac{df\left(  g\right)  }{dg}\frac{d^{2}g\left(  x\right)
}{dx^{2}}\right\} \nonumber\\
&  =\frac{d^{3}f\left(  g\right)  }{dg^{3}}\left(  \frac{dg\left(  x\right)
}{dx}\right)  ^{3}+3\frac{d^{2}f\left(  g\right)  }{dg^{2}}\frac{dg\left(
x\right)  }{dx}\frac{d^{2}g\left(  x\right)  }{dx^{2}}+\frac{df\left(
g\right)  }{dg}\frac{d^{3}g\left(  x\right)  }{dx^{3}}\nonumber\\
&  =\frac{d^{3}f\left(  g\right)  }{dg^{3}}\left(  \frac{dg\left(  x\right)
}{dx}\right)  ^{3}-3\left\{  \left(  \frac{dg\left(  x\right)  }{dx}\right)
^{-1}\frac{d^{2}g\left(  x\right)  }{dx^{2}}\left(  \frac{dg\left(  x\right)
}{dx}\right)  ^{-2}\right\}  \frac{dg\left(  x\right)  }{dx}\frac
{d^{2}g\left(  x\right)  }{dx^{2}}\nonumber\\
&  +\left(  \frac{dg\left(  x\right)  }{dx}\right)  ^{-1}\frac{d^{3}g\left(
x\right)  }{dx^{3}}\nonumber\\
&  \Rightarrow\frac{d^{3}f\left(  g\right)  }{dg^{3}}=\left(  \frac{dg\left(
x\right)  }{dx}\right)  ^{-3}\left\{  3\left\{  \left(  \frac{dg\left(
x\right)  }{dx}\right)  ^{-1}\frac{d^{2}g\left(  x\right)  }{dx^{2}}\left(
\frac{dg\left(  x\right)  }{dx}\right)  ^{-1}\right\}  \frac{d^{2}g\left(
x\right)  }{dx^{2}}\right. \nonumber\\
&  \qquad\qquad\qquad\qquad\qquad\qquad\qquad\qquad\qquad\qquad\qquad
\qquad\left.  -\left(  \frac{dg\left(  x\right)  }{dx}\right)  ^{-1}%
\frac{d^{3}g\left(  x\right)  }{dx^{3}}\right\}
\end{align}



\end{document}